%!TEX root = ../template.tex
%%%%%%%%%%%%%%%%%%%%%%%%%%%%%%%%%%%%%%%%%%%%%%%%%%%%%%%%%%%%%%%%%%%%
%% Config/4_files.tex
%% NOVA thesis configuration file
%%%%%%%%%%%%%%%%%%%%%%%%%%%%%%%%%%%%%%%%%%%%%%%%%%%%%%%%%%%%%%%%%%%%

\typeout{NT FILE Config/4_files.tex}%

%=============================================================================
% THE FILES
%=============================================================================

%-------------------------------------------------------------
%% Except for the bibliography, all the other FILE NAMES below
%%     inside braces must correspond to a file with extension
%%     ".tex" and located in the "Chapters" folder
%-------------------------------------------------------------

%-------------------------------------------------------------
% BibTeX bibliography files and customization.
%.............................................................
%
% May be used multiple times with a single file name each time
%-------------------------------------------------------------
\ntaddfile{bib}{bibliography.bib}
%\ntaddfile{bib}{anotherbibfile.bib}



%-------------------------------------------------------------
% File with the dedication text.
%.............................................................
%
% Ignored if 'status=working'
%-------------------------------------------------------------
\ntaddfile{dedication}{dedication}



%-------------------------------------------------------------
% File with the acknowledgments text.
%.............................................................
%
% Ignored if 'status=working'
%-------------------------------------------------------------
\ntaddfile{acknowledgements}{acknowledgements}



%-------------------------------------------------------------
% File with the statement text.
%.............................................................
%
% Ignored if 'status=working'
%-------------------------------------------------------------
% \ntaddfile{statement}[en]{univ-school/statement-en}
% \ntaddfile{statement}[pt]{univ-school/statement-pt}



%-------------------------------------------------------------
% File with the qute text.
%.............................................................
%
% Ignored if 'status=working'
%-------------------------------------------------------------
\ntaddfile{quote}{quote}



%-------------------------------------------------------------
% Abstracts
%
%     \ntaddfile{abstract}[<lang>]{<filename>}
%.............................................................
%
% <lang>: -> cat (catalan), cz (czech), de (german), dk (danish), en (english), 
%            es (spanish), fr (french), gr (greek), it (italian), nl (dutch), 
%            pl (polish), pt (portuguese), sk (slovak), uk (ukrainian)
%            zhs (chinese simplified), zht (chinese traditional)
%
% <filename>: -> must be '1-Frontmatter/filename.tex'
%
% Specifically for ULISBOA-FMV, the extended asbtract
%       \ntaddfile{abstract}[pt-ext]{abstract-pt-ext}
%-------------------------------------------------------------
\ntaddfile{abstract}[pt]{abstract-pt}  % Abstract in Portuguese
\ntaddfile{abstract}[en]{abstract-en}  % Abstract in English
% \ntaddfile{abstract}[de]{abstract-de}  % Abstract in German
% \ntaddfile{abstract}[es]{abstract-es}  % Abstract in Spanish
% \ntaddfile{abstract}[fr]{abstract-fr}  % Abstract in French
% \ntaddfile{abstract}[it]{abstract-it}  % Abstract in Italian
% \ntaddfile{abstract}[gr]{abstract-gr}  % Abstract in Greek
% \ntaddfile{abstract}[uk]{abstract-uk}  % Abstract in Ukrainian
% \ntaddfile{abstract}[zhs]{abstract-zhs}  % Abstract in Chinese (Simplified)
% \ntaddfile{abstract}[zht]{abstract-zht}  % Abstract in Chinese (Traditional)
% % \ntaddfile{abstract}[pt-ext]{abstract-pt-ext}  % Extended abstract in Portuguese (e.g., for ulisboa/fmv)



%-------------------------------------------------------------
% Glossaies (glossary, acronyms, symbols, etc)
%
%     \ntaddfile{glossaries}[<type>]{<filename>}
%.............................................................
%
% <filename>: -> must be '1-Frontmatter/filename.tex'
%
% <type>:  glossary -> Glossay (longer dictionary like entries)
%           acronym -> Acronyms and their meaning
%           symbols -> Symbols and their meaning
%-------------------------------------------------------------
\ntaddfile{glossaries}[glossary]{glossary}
\ntaddfile{glossaries}[acronym]{acronyms}
\ntaddfile{glossaries}[symbols]{symbols}
% \ntaddfile{glossaries}[symbols]{chemical}



%-------------------------------------------------------------
% Main text - document chapters
%
\ntaddfile{chapter}{chapter-1-introducao}
\ntaddfile{chapter}{chapter-2-enquadramento}
\ntaddfile{chapter}{chapter-3-revisao}
\ntaddfile{chapter}{chapter-4-metodos}
\ntaddfile{chapter}{chapter-5-plataforma}
\ntaddfile{chapter}{chapter-6-diagnostico}
\ntaddfile{chapter}{chapter-7-discussao}
\ntaddfile{chapter}{chapter-8-conclusoes}



%-------------------------------------------------------------
% Apppendices
%
\ntaddfile{appendix}{appendix-A-protocolo-sr}
\ntaddfile{appendix}{appendix-B-validacao}
\ntaddfile{appendix}{appendix-C-dicionario}
\ntaddfile{appendix}{appendix-D-tabelas}
\ntaddfile{appendix}{appendix-E-manual}
\ntaddfile{appendix}{appendix-F-adversarial}
\ntaddfile{appendix}{appendix-G-matematica-rascunho}



%-------------------------------------------------------------
% Apppendices
%
% (sem annexs)



%-------------------------------------------------------------
% Covers
%
%     \ntaddfile{cover}[<cover-id>]{<filename>}
%.............................................................
%
% <filename>:   -> must be '1-Frontmatter/filename.tex'
%
% <cover-id>: 1 -> The main cover
%             2 -> The front page
%             N -> The back cover
%         spine -> The back cover
%
% The given file will be used to render the cover <cover-id>
% Note: if you want to skip prnting the frontpage, do not
%       redefine the frontpage! Instead use the command
%       \ntsetup{print/frontpage=false}
%-------------------------------------------------------------
% \ntaddfile{cover}[1]{cover-page}


